\section{Results}
\label{sec:eval}

\subsection{How we tested}
\label{sec:setup}

We did not choose shapes by hand. Take a tensor with
$N = \prod_{i=1}^{m} p_i^{a_i}$ values and rank $r$. Every valid shape comes
from splitting each prime's exponent across the $r$ dimensions, so that
$\sum_{j=1}^{r} b_{ij} = a_i$:

\begin{equation}
s = \left( \prod_{i=1}^{m} p_i^{b_{i1}},\ \prod_{i=1}^{m} p_i^{b_{i2}},\
           \ldots,\ \prod_{i=1}^{m} p_i^{b_{ir}} \right).
\label{eq:shapes}
\end{equation}

We draw from that whole set. So nobody picked the cases. We drop shapes with an
axis of size one. Such an axis can move anywhere without touching a byte.

The main test has \resCases{} cases at ranks \resRankLo{} to \resRankHi{}.
Every tensor has \resElements{} float32 values, which is \resTensorMB{}\,MB. We
ran them on a \resDevice{}. Each shape is paired with a permutation that leaves
a group of axes at the back alone. That is what our method is built for.

We ran a second test with \hardCases{} permutations picked at random. That
covers what our method is not built for. Section~\ref{sec:limits} reports it.

We compare against PyTorch \texttt{.contiguous()}. We also compare against our
rebuilds of the three published in-place methods.
Section~\ref{sec:related} lists them. We time every method with CUDA events
after one warmup run. We check every answer against
\texttt{permute(*perm).contiguous()} before it counts.

Our method finished and gave the right answer on every case in both tests. It
turned nothing away.

% Table (tab:methods). Body and caption macros: _generated/methods-table.tex
% and .defs.tex, written by pipeline/tables/gen_table_methods.py from
% raw_data/dpost_sweep/.
%
% table*, not table. Five columns -- a method name plus two byte-counts wide
% enough to carry a unit -- have a natural width well past the 241pt of a
% single USENIX column. tabular* does NOT shrink to fit: it only distributes
% slack, so when the content is already too wide it runs past the column and
% over the neighbouring text WITHOUT raising an Overfull warning. That is why
% this has to be caught by eye, and why the full measure is the fix rather
% than \small-ing the type further.
\begin{table*}[!t]
\centering
\caption{\capheadline{Every method, every case, no routing.}
\methTabCases{} shape and permutation pairs at ranks \resRankLo{}--\resRankHi{},
\resTensorMB{}\,MB per tensor, on a \resDevice{}. We run our method on all of
them, including the ones our own rule would turn away.}
\label{tab:methods}
\small
\begin{minipage}{0.72\linewidth}
\centering
% GENERATED by pipeline/tables/gen_table_methods.py -- do not edit by hand.
% Regenerate after any change to raw_data/dpost_sweep/.
\begin{tabular}{lrrrr}
\toprule
Method & Median & Median aux & Peak aux & Box wins \\
\midrule
PyTorch \texttt{.contiguous()} & 7.70 & 1,024\,MB & 1,024\,MB & 32/87 \\
Box transpose (ours) & \textbf{10.21} & \textbf{144\,B} & \textbf{10.6\,KB} & -- \\
ITTPD~\cite{cheng2025ittpd} & 80.49 & 0 & 0 & 87/87 \\
Gomez-Luna~\cite{gomezluna2016inplace} & 82.91 & 16\,B & 512\,B & 87/87 \\
Gustavson~\cite{gustavson2012cycleleader} & 97.39 & 48\,MB & 128\,MB & 79/87 \\
\bottomrule
\end{tabular}

\vspace{1.5pt}
\parbox{\linewidth}{\fontsize{6}{6.8}\selectfont\raggedright Middle-of-the-range values over every case, with no routing: we run our method on all of them, including the ones our own rule would turn away. \textbf{Box wins} counts cases where box transpose is faster than that method. Against \texttt{.contiguous()} that count is low here precisely because no routing is applied; Section~\ref{sec:gate} reports it under the gate. ITTPD's memory figure comes from its own planner, not from watching the allocator. On these shapes every block fits in fast on-chip memory, so it never asks for anything global.}
%
\end{minipage}
\end{table*}


\subsection{Against everyone}

Table~\ref{tab:methods} shows all \methTabCases{} cases. We used no routing
here. We ran our method even on the cases our own rule says to skip.

In the middle of the range we are \methTabIttpdX{}$\times$ faster than ITTPD,
\methTabGlX{}$\times$ faster than Gomez-Luna, and \methTabGvX{}$\times$ faster
than Gustavson. Case by case, we are faster on \resWinIttpd{} of
\resWinIttpdOf{}, on \resWinGl{} of \resWinGlOf{}, and on \resWinGv{} of
\resWinGvOf{}.

The memory column matters too. \texttt{.contiguous()} needs a second tensor
every time. We need \resBoxAux{} bytes. Gustavson needs \resGvAuxMB{}\,MB.
ITTPD needs none. That is real, not a mistake. On these shapes its blocks fit
in fast on-chip memory, so it never asks for global memory. Against ITTPD we
win on speed, not on memory. We do not claim otherwise.

% Table (tab:tiers). Body and caption macros: _generated/tiers-table.tex and
% .defs.tex, written by pipeline/tables/gen_table_tiers.py from both sweeps.
%
% table*, not table: seven columns of data do not fit a 241pt column.
\begin{table*}[!t]
\centering
\caption{\capheadline{Three cases, and who wins each.}
Both sweeps together, split by the two numbers from
Section~\ref{sec:method}. $\tau$ is how long a run has to be before the memory
system reads it efficiently, \gateMinDpost{} values. The first row is where we
beat the copy. The last is where our method is the wrong tool.}
\label{tab:tiers}
\small
% GENERATED by pipeline/tables/gen_table_tiers.py -- do not edit by hand.
% Regenerate after any change to raw_data/dpost_sweep/ or raw_data/hard_perms/.
\begin{tabular*}{\linewidth}{@{\extracolsep{\fill}}lrrrrrr@{}}
\toprule
 & & \multicolumn{3}{c}{Median ms} & \multicolumn{2}{c}{Box faster than} \\
\cmidrule(lr){3-5}\cmidrule(lr){6-7}
Regime & Cases & Box & Copy & Best other & Copy & Best other \\
\midrule
$D_{post}\!\ge\!\tau$, one small square swap & 22 & 4.03 & 7.68 & 8.16 & 22/22 & 22/22 \\
$D_{post}\!\ge\!\tau$, otherwise & 69 & 12.48 & 7.69 & 22.37 & 11/69 & 61/69 \\
$D_{post}\!<\!\tau$ & 11 & 228.12 & 7.48 & 74.28 & 0/11 & 0/11 \\
\bottomrule
\end{tabular*}

\vspace{1.5pt}
\parbox{\linewidth}{\fontsize{6}{6.8}\selectfont\raggedright \textbf{Copy} is PyTorch \texttt{.contiguous()}. \textbf{Best other} is whichever in-place method was fastest on that case. Both tests are worked out from the shape before any data is touched. The last row is where our method falls down: a permutation that moves the innermost axis forces a swap with $D_{post}=1$, each value we move is then a lone number, and the published in-place methods are faster. We report this rather than route around it.}
%
\end{table*}


\subsection{Three cases}
\label{sec:gate}

Table~\ref{tab:tiers} splits both tests by the two numbers from
Section~\ref{sec:method}. We work out both from the shape before touching any
data.

The back group is large and the swap is one small square block. Here we are
faster than everyone, including the copy. We won \tierBeatsCases{} of
\tierBeatsCases{} cases, at \tierBeatsX{}$\times$ in the middle of the range.
In the main test alone, our rule fires on \gateShare{}\% of shapes and wins
\gateWins{} of \gateCases{}. That is \gateMedianX{}$\times$ in the middle and
\gateBestX{}$\times$ at best.

The back group is large but the swap does not qualify. Here we lose to the
copy. We are still the best in-place choice, on \tierInplaceCases{} cases.

The back group is too small. Here we are the worst. Section~\ref{sec:limits}
covers it.

\subsection{The prediction was right}

We did not fit the boundary to the results. Equation~(\ref{eq:moved}) says a
swap of a $\modelMinMovedRows \times \modelMinMovedCols$ block moves
$\modelMinMoved$ of the tensor. So it should take $\modelMinMoved$ of the
copy's time. From rank \resRankLo{} to rank \resRankHi{}, the measured times
land on that value. The error inside the region is about $\modelErrIn\%$.

The speedup comes from arithmetic, not from tuning. It does not depend on how
well we wrote the kernel. It does not depend on which GPU we used. It depends
only on how much of the tensor was already in the right place.
